% ID: SEL-D0066-OB-E02
% ESTADO: cotejada. Segunda pasada visual y borrador compilado revisados.
% Fuente: Selectividad/sel_2010_matematicas_aplicadas/mat_apli/mat_apli_sept/mat_apli_sept_rb/EXAMEN-ANDALUCÍA-2.pdf
% SHA-256: 24c494a5d013dbc7f5adaac9896dee2d79f22540518359c67eac29945a26d4f0
% Página: 2; opción B; ejercicio original 2.

El gerente de una empresa sabe que los beneficios de la misma, $f(x)$, dependen de la inversión, $x$, según la función $f(x)=-x^2+11x-10$.

($x$ es la cantidad invertida, en millones de euros).
\begin{enumerate}[label=\alph*)]
\item \textbf{(0.75 puntos)} Determine los valores de la inversión para los que la función beneficio es no negativa.
\item \textbf{(1 punto)} Halle el valor de la inversión para el cual el beneficio es máximo. \textquestiondown{}A cuánto asciende éste?
\item \textbf{(0.75 puntos)} \textquestiondown{}Entre qué valores ha de estar comprendida la inversión para que el beneficio sea creciente, sabiendo que éste es no negativo?
\end{enumerate}
